\subsection{The quintic: a certified finite-descent target}

\begin{statement}{A.50}{The quintic genus-two target and its exact modular input}{Known}
Let
\[
 C:\ y^2+y=x^5,\qquad J=\operatorname{Jac}(C),\qquad
 F=\Q(\sqrt5),\qquad K=\Q(\zeta_5).
\]
The point at infinity of $C$ is rational. The abelian surface $J$ has
complex multiplication by $K$, with all endomorphisms defined over $K$,
and real multiplication by $F$ over $F$. Over $F$ it is associated with
the two conjugate parallel-weight-two Hilbert eigenforms of trivial
central character and level $\mathfrak p^3$, where
$\mathfrak p=(\sqrt5)$. Its conductor over $\Q$ is $5^5$, and
\[
 L(J/F,s)=L(J/\Q,s)^2.
\]
Consequently the established Hilbert modular rank-zero theorem implies
that $J(F)$ and $\Sha(J/F)$ are finite as soon as $L(J/\Q,1)\ne0$.
\end{statement}
\begin{proof}[Verification of the stated input and all its hypotheses]
The automorphism $(x,y)\mapsto(\zeta_5x,y)$ acts on $J$. On the
holomorphic differentials $dx/(2y+1)$ and $x\,dx/(2y+1)$ its eigenvalues
are $\zeta_5$ and $\zeta_5^2$. Thus the CM type is primitive: the only
proper nontrivial subfield of $K$ is real. In particular $J$ is absolutely
simple and $\operatorname{End}^0_{\overline\Q}(J)=K$. Galois acts on this
algebra through its usual action on $\zeta_5$, so
\[
 \operatorname{End}^0_\Q(J)=\Q,\qquad
 \operatorname{End}^0_F(J)=F.
\]
The Rosati involution is complex conjugation, and the real multiplication
is self-adjoint. This explains both why the direct modular-curve theorem
over $\Q$ cannot be used and why a Hilbert modular argument over $F$ can.

The precise modular identification is proved in
\citet[p.~439, the paragraph following Table~3]{AMDarmon2000}: the Hecke
characters have conductor $(1-\zeta_5)^2$, and the resulting two forms
over $F$ have level $\mathfrak p^3$. These are newforms, since the
lower-level spaces in Table~2 are zero. Their coefficient field is $F$;
for example the displayed eigenvalue at a prime of norm eleven is
$(-1-5\sqrt5)/2$. The CM construction is the modularity input here,
not an inference from a finite match of Fourier coefficients.

For clarity, the conductor and $L$-function comparison can also be
recovered directly from these characters. After extending the coefficient
field, the four-dimensional Tate representation over $\Q$ is
\[
 V_\ell J\otimes_{\Q_\ell}\overline\Q_\ell\simeq
 \operatorname{Ind}_{G_K}^{G_\Q}\psi_\ell.
\]
Indeed its four distinct eigenspaces for $\zeta_5$ are permuted
transitively by $G_\Q$, and the stabilizer of one is $G_K$.
The conductor formula for induction gives
\[
 N(J/\Q)=|\operatorname{disc}K|\,
     \operatorname{Norm}_{K/\Q}(1-\zeta_5)^2
 =5^3\cdot5^2=5^5.
\]
Since these characters are ramified at the unique prime above five,
the inertia invariants there are zero and the bad Euler factor is one.
Let $\chi_5$ be the quadratic character of $F/\Q$. It is trivial on
$G_K$, so tensoring this induced representation by $\chi_5$ leaves it
unchanged. Artin formalism therefore gives, including the bad factors,
\[
 L(J/F,s)=L(J/\Q,s)L(J/\Q\otimes\chi_5,s)
          =L(J/\Q,s)^2.
\]
On the other hand modularity over $F$ gives
$L(J/F,s)=L(f,s)L(f^\tau,s)$, in arithmetic normalization with center
$s=1$. Thus nonvanishing of the right-hand square implies nonvanishing
of each Hilbert factor separately.

The arithmetic theorem used is the central-rank-at-most-one theorem
stated by \citet[Theorem~4.3.2, p.~588]{AMZhang2010}: for a parallel
weight-two cuspidal Hilbert representation, the corresponding abelian
variety has the predicted rank and finite Tate--Shafarevich group if
the base field has odd degree or some finite local component is not
principal series. We verify the second alternative. The central character
is trivial, and the conductor exponent at $\mathfrak p$ is three.
A principal series with trivial central character has characters inverse
to one another on units, so its conductor exponent is twice the conductor
of either character. It cannot have exponent three. Hence the local
component at $\mathfrak p$ is not principal series. Equivalently, it
transfers to the quaternion algebra ramified at $\mathfrak p$ and at
one of the two real places, so the required Shimura-curve realization
exists. All field, weight, character, local, and analytic hypotheses of
the stated theorem have now been checked.

The source chain is recorded precisely. The printed reference for this
theorem in Zhang's article is Tian--Zhang, ``in preparation.'' We use the
explicitly stated published theorem as a \emph{Known input}; we do not
claim to have opened that unpublished work. Its Euler-system formulation
is Nekov\'a\v r's Theorem~3.2, and the Gross--Zagier formula in the
coprime-discriminant case is \citet[Theorem~6.1, p.~203]{AMZhang2004}.
In particular, the older exponent-one hypothesis in
\citet[Theorem~A]{AMZhang2001} is not silently replaced by exponent three.
For the Euler-system CM exclusion, one chooses the auxiliary CM quadratic
extension of $F$ unramified and inert at $\mathfrak p$; it differs from
$K/F$, which is ramified there. The analytic nonvanishing theorem for
such auxiliary quadratic twists is part of the published rank-zero
input just specified. No finiteness claim is obtained merely from real
multiplication or from algebraic rank zero.
\end{proof}

\begin{statement}{A.51}{A rigorous degree-four central-value certificate}{Proved}
For $J=\operatorname{Jac}(y^2+y=x^5)$, the root number is $+1$ and
\[
 L(J/\Q,1)\in
 [\,0.71831362842056477051660347805
       \mathbin{\pm}2.72\cdot10^{-30}\,].
\]
The certificate uses exact Euler factors, $4096$ Dirichlet coefficients,
and an explicit bound for the entire omitted tail.
\end{statement}
\begin{proof}
Put $N=3125$, $c=4\pi^2/\sqrt N$, and
$L(s)=\sum_{n\ge1}a_nn^{-s}$. The CM Hecke-character description in A.50
gives an entire completed function
\[
 \Lambda(s)=c^{-s}\Gamma(s)^2L(s),\qquad
 \Lambda(s)=w\Lambda(2-s),\quad w\in\{1,-1\}.
\]
The omitted constant factor four in the usual product of two complex
gamma factors does not change the functional equation. Set
\[
 \Theta(t)=2\sum_{n\ge1}a_nK_0(2\sqrt{cnt}),\quad t>0.
\]
The Mellin identity
$\int_0^\infty2K_0(2\sqrt x)x^{s-1}dx=\Gamma(s)^2$
gives
$\Lambda(s)=\int_0^\infty\Theta(t)t^{s-1}dt$ initially in its
half-plane of absolute convergence. Mellin inversion and the functional
equation imply
$\Theta(t)=wt^{-2}\Theta(1/t)$. In particular a certified nonzero
value of $\Theta(1)$ determines $w=+1$ without guessing a numerical
root number. Once that sign is known, splitting the Mellin integral
at one and using
$\int_1^\infty K_0(2\sqrt{cnt})dt
 =K_1(2\sqrt{cn})/\sqrt{cn}$ gives
\[
 L(1)=4\sqrt c\sum_{n\ge1}\frac{a_n}{\sqrt n}
           K_1(2\sqrt{cn}).
\]

Here is a sufficient elementary tail estimate. The Weil bounds on each
degree-four Euler factor imply
$|a_n|\le d_4(n)\sqrt n\le n^{7/2}$; the last inequality follows by
counting the first three entries of an ordered factorization into four
positive factors. For $x\ge2$, the integral representation
\[
 K_\nu(x)=\int_0^\infty e^{-x\cosh u}\cosh(\nu u)du,
 \qquad \nu=0,1,
\]
together with $\cosh u\ge1+u^2/2$ and
$\log\cosh u\le u^2/2$, gives
$K_\nu(x)\le\sqrt{\pi/x}\,e^{-x}$.
This integral representation is also
\citet[equation~10.32.9]{AMDLMF2026}.
Write $k=2\sqrt c$. For $M$ with $k\sqrt M>8$, the function
$x^4e^{-k\sqrt x}$ is decreasing for $x\ge M$, and therefore
\[
 \sum_{n>M}n^4e^{-k\sqrt n}
 \le\int_M^\infty x^4e^{-k\sqrt x}dx
 =\frac{2\Gamma(10,k\sqrt M)}{k^{10}}=:I_M.
\]
Thus the errors after truncating the theta and central-value series
are at most
\[
 T_\Theta=2\sqrt{\pi/k}\,I_M,\qquad
 T_L=4\sqrt c\sqrt{\pi/k}\,I_M.
\]
No unevaluated improper integral remains in this bound, since
$\Gamma(10,z)=9!e^{-z}\sum_{j=0}^9z^j/j!$.

The executed script \texttt{a\_main\_quintic.py} computes every good
Euler polynomial for $p\le4096$ directly over the finite field using
\texttt{hyperellcharpoly}; the factor at five is one. It expands their
reciprocals independently in Python and checks every coefficient against
PARI's exact \texttt{lfunan} output. The conductor is also checked by
\texttt{genus2red}, which returns $3125$ and type VIII-1 at five.
The model is smooth at two, so the documented uncertainty of the
genus-two conductor routine at bad even reduction does not occur.
For an explicit good-reduction check, the derivative in $y$ is one
modulo two on the affine chart; the chart at infinity is
$v^2+u^3v=u$ with $u=1/x$, $v=y/x^3$, whose derivative in $u$
at $(0,0)$ is $-1$. The precise routines and their caveats are documented
by \citet{AMPARIHyperell2026}.
All arithmetic input factors and coefficients are saved in
\texttt{a\_main\_quintic\_data.json.gz} and the raw GP output is retained.

With $192$-bit Arb arithmetic and $M=4096$, the outward-rounded
calculations give
\[
 \Theta(1)\in[0.3300364476\mathbin{\pm}3.40\cdot10^{-11}],
 \quad T_\Theta<1.231\cdot10^{-30},
 \quad T_L<2.068\cdot10^{-30}.
\]
The first interval excludes zero and proves $w=+1$. The central-value
interval, after adding the entire tail, is the one displayed in the
statement and is positive. No tabulated analytic rank, approximate zero
test, guessed sign, or uncertified \texttt{lfun} value enters this proof.
\end{proof}

\begin{statement}{A.52}{Every maximal-prosolvable Selmer section of the quintic is cuspidal}{Proved}
Let $U_5/\Q$ be the plus Fermat curve
$X^5+Y^5+Z^5=0$ with all coordinate cusps removed. Every Selmer section
of its geometrically maximal prosolvable arithmetic fundamental group
is in the image of a rational cusp decomposition group. This conclusion
uses the explicitly identified Known Hilbert modular input of A.50.
\end{statement}
\begin{proof}
By A.50--A.51, $J(F)$ and $\Sha(J/F)$ are finite. It follows that
$J(\Q)$ is finite. The restriction map on Tate--Shafarevich groups has
finite image, and its kernel is contained, by inflation--restriction, in
\[
 H^1(\Gal(F/\Q),J(F)),
\]
which is finite. Hence $\Sha(J/\Q)$ is finite too. In particular
there is no nonzero divisible part. With infinity as the degree-one
base divisor, the rank-zero finite-abelian-descent criterion gives
\[
 C(\mathbb A_\Q)_\bullet^{\mathrm{f\text{-}ab}}=C(\Q).
\]
This is the Jacobian criterion used in
\citet[section~6]{AMStoll2007}; all its arithmetic hypotheses have just
been supplied.

The rational points can be determined without a rational-point table.
The exact Euler polynomials at three and eleven give
\[
 \#J(\F_3)=10,\qquad \#J(\F_{11})=125.
\]
Prime-to-residue-characteristic torsion injects under good reduction.
For primes other than three and eleven the torsion order is bounded
by $\gcd(10,125)=5$; the eleven-adic bound excludes three-primary
torsion, and the three-adic bound excludes eleven-primary torsion.
Thus $\#J(\Q)$ divides five. If $R=(0,0)$ and $R'=(0,-1)$, then
\[
 \operatorname{div}(y)=5R-5\infty,\qquad
 \operatorname{div}(x)=R+R'-2\infty.
\]
The class $D=[R-\infty]$ is nonzero, since a principal divisor
$R-\infty$ would give a degree-one map to $\mathbf P^1$.
Consequently $J(\Q)=\langle D\rangle\simeq\mathbf Z/5$.

The Abel--Jacobi image of $C(\Q)$ can therefore consist only of
$0,\pm D,\pm2D$. The first three are $\infty,R,R'$. If
$[P-\infty]=2D$, then $P+\infty\sim2R$. The point $R$ is not a
Weierstrass point: the hyperelliptic involution sends it to the distinct
point $R'$. On a genus-two curve the only degree-two linear system of
positive dimension is the hyperelliptic system. Thus $h^0(2R)=1$ and
$P+\infty=2R$, an impossibility. The same argument with $R'$ excludes
$-2D$. This proves
\[
 C(\Q)=\{\infty,(0,0),(0,-1)\}.
\]

On the plus Fermat model the quotient map is
\[
  x=-\frac{XY}{Z^2},\qquad y=\frac{X^5}{Z^5}.
\]
Indeed $y^2+y=-X^5Y^5/Z^{10}=x^5$. It extends to the smooth proper
curves, and the inverse image of the displayed three rational points
is precisely the coordinate-cusp divisor. On the classical model
$X^5+Y^5=Z^5$ the formula is instead
$y=-X^5/Z^5$; the conversion is $Z_{\rm classical}=-Z_{\rm plus}$.
This sign distinction is necessary.

Apply the finite-support theorem B.1 to a maximal-prosolvable Selmer
section and this map. Its image on $C$ gives an adelic point surviving
finite abelian descent and therefore one of the three points above.
All local supports on the Fermat curve lie in the finite cusp fibre.
B.1 uses only solvable neighbourhoods and finite abelian refinements,
which remain visible in the maximal prosolvable quotient. It produces
a rational point of that cusp fibre and a section of its decomposition
group mapping to the given section. This proves the stated quotient
claim; no full-profinite lift is asserted or needed.
\end{proof}

\begin{statement}{A.53}{The entire quintic Jacobian is a cube of the certified target}{Proved}
Let $X_5$ be the smooth projective plus Fermat quintic over $\Q$ and
$J=\operatorname{Jac}(y^2+y=x^5)$. There is a $\Q$-isogeny
\[
 \operatorname{Jac}(X_5)\longrightarrow J^3.
\]
Consequently $\operatorname{Jac}(X_5)(\Q)$ and
$\Sha(\operatorname{Jac}(X_5)/\Q)$ are finite, with the Known arithmetic
input stated in A.50. This supplies compact abelian rigidity for the
quintic as well as the maximal-prosolvable conclusion of A.52.
\end{statement}
\begin{proof}
In addition to the map in A.52, cyclically permute the coordinates to
obtain three maps $f_i:X_5\to C$:
\[
 \left(-\frac{XY}{Z^2},\frac{X^5}{Z^5}\right),\quad
 \left(-\frac{YZ}{X^2},\frac{Y^5}{X^5}\right),\quad
 \left(-\frac{ZX}{Y^2},\frac{Z^5}{Y^5}\right).
\]
Each satisfies $v^2+v=u^5$, and each rational map extends to the proper
smooth curves. On the chart $Z=1$, write
$\omega=dX/Y^4$ and use
$dY=-X^4Y^{-4}dX$. The pullbacks of the basis
$du/(2v+1),\ u\,du/(2v+1)$ of $H^0(C,\Omega^1)$ are
\[
 \begin{array}{c|cc}
    &du/(2v+1)&u\,du/(2v+1)\\ \hline
 f_1&\omega&-XY\omega\\
 f_2&X^2\omega&-Y\omega\\
 f_3&Y^2\omega&-X\omega.
 \end{array}
\]
For example,
$d(-XY)=(X^5-Y^5)\omega$ and
$2X^5+1=X^5-Y^5$, which gives the first entry. The others follow by
the same differentiation and the equation $X^5+Y^5+1=0$.

Adjunction for a smooth plane quintic identifies its six-dimensional
space of regular differentials with
\[
 \langle 1,X,Y,X^2,XY,Y^2\rangle\,\omega.
\]
The displayed six pullbacks are exactly this basis up to permutation
and signs. Thus the homomorphism
$(f_{1*},f_{2*},f_{3*}):\operatorname{Jac}(X_5)\to J^3$
induces an isomorphism on cotangent spaces. The source and target have
the same dimension six, so this homomorphism is a $\Q$-isogeny.
The executed exact symbolic certificate
\texttt{a\_main\_quintic\_maps.py} reduces all three curve identities
and all six differential identities modulo $X^5+Y^5+1$, and computes
the coefficient matrix determinant to be one. This determinant proves
independence; it makes no claim that the isogeny has degree one.

Rank is invariant under isogeny, and finiteness of the
Tate--Shafarevich group is also invariant: the kernel of the induced
map on $\Sha$ is controlled by a finite isogeny Selmer group, and a
dual isogeny supplies the reverse direction. Hence the finiteness
established for $J$ in A.52 passes to $\operatorname{Jac}(X_5)$.
With a rational coordinate cusp as base point, Stoll's rank-zero
Jacobian criterion applies to $X_5$ itself. Its rational points are
all coordinate cusps, since any rational point maps under $f_1$ to
one of the three points in A.52, whose complete fibre is the
coordinate-cusp divisor. The resulting first-derived-quotient
argument is given in B.26.
\end{proof}

\begin{statement}{A.54}{The local quintic Jacobians have no prime-to-five torsion}{Proved}
For $J=\operatorname{Jac}(y^2+y=x^5)$ and
$A=\operatorname{Jac}(X_5)$, both $J(\Q_5)$ and $A(\Q_5)$ are pro-$5$
groups. In particular the local Kummer maps
\[
 J(\Q_5)\longrightarrow H^1(\Q_5,T_5J),\qquad
 A(\Q_5)\longrightarrow H^1(\Q_5,T_5A)
\]
are injective. The assertion uses the exact CM conductor in A.50;
it does not require its rank or Tate--Shafarevich conclusion.
\end{statement}
\begin{proof}
Put $K_5=\Q_5(\zeta_5)$ and $\lambda=1-\zeta_5$.
Restrict the induced CM representation from A.50 to $G_{K_5}$.
Its four characters all have conductor exactly $\lambda^2$:
they are the four conjugates of the character of that conductor
identified by \citet[p.~439]{AMDarmon2000}, and $K$ has a unique
prime above five. For each coefficient prime $\ell\ne5$, local
reciprocity identifies the finite inertia character with the finite
local component of the corresponding CM Hecke character. Thus it is
trivial on $1+\lambda^2\mathcal O_{K_5}$ and nontrivial on
\[
 (1+\lambda\mathcal O_{K_5})/
 (1+\lambda^2\mathcal O_{K_5})\simeq(\F_5,+).
\]
Choose a wild-inertia element $\tau$ whose image generates this
quotient. Every one of the four eigenvalues of $\tau$ is a primitive
fifth root of unity, so
$\Phi_5(\tau)=0$ on $V_\ell J$, where
$\Phi_5(T)=T^4+T^3+T^2+T+1$.

This polynomial identity holds on every stable integral Tate lattice.
The integral identity
\[
 (T-1)(T^3+2T^2+3T+4)=\Phi_5(T)-5
\]
then shows that $\tau-1$ is invertible on $T_\ell J$, and hence on
$J[\ell^\infty]$, whenever $\ell\ne5$. In particular
$J(\Q_5)[\ell^\infty]=0$ for every such $\ell$, including two.
The $\Q$-isogeny of A.53 gives
$V_\ell A\simeq(V_\ell J)^3$. Therefore the same polynomial identity
holds on $V_\ell A$ and then on its own integral Tate lattice, giving
$A(\Q_5)[\ell^\infty]=0$ too. This argument needs no restriction on
the degree of the isogeny.

The group of points of an abelian variety over $\Q_5$ is compact and
contains an open subgroup isomorphic to a finite power of $\Z_5$.
Its prime-to-five Sylow subgroups are therefore finite torsion groups;
the preceding vanishing shows they are zero. Both local point groups
are consequently pro-$5$, and as finitely generated $\Z_5$-modules
with finite $5$-primary torsion they satisfy
$\bigcap_{r\ge1}5^rJ(\Q_5)=0$ and
$\bigcap_{r\ge1}5^rA(\Q_5)=0$. These are exactly the kernels of the
two displayed Kummer maps.

One can also retain the degree information in the explicit geometry.
Each map $f_i$ in A.53 has degree five. Put
$\Phi=(f_{1*},f_{2*},f_{3*})$ and
$\Psi=f_1^*+f_2^*+f_3^*$.
Over $\overline\Q$, use the action
$(X,Y)\mapsto(\zeta_5^aX,\zeta_5^bY)$ on $Z=1$.
The two holomorphic character vectors for the three quotient factors
are respectively
\[
 \{(1,1),(2,2)\},\quad
 \{(3,1),(1,2)\},\quad
 \{(1,3),(2,1)\}.
\]
They lie on distinct lines in $\F_5^2$. The respective order-five
quotient groups have generators $(1,4),(1,2),(2,1)$.
Their trace operators act as five on their own factor and zero on
each of the other two factors. Hence
$f_{i*}f_j^*=0$ for $i\ne j$, whereas
$f_{i*}f_i^*=[5]$. This proves
$\Phi\Psi=[5]$ on $J^3$ and, since $\Phi$ is an isogeny,
$\Psi\Phi=[5]$ on $A$. Thus both kernels are killed by five.
The executed program \texttt{a\_main\_quintic\_inertia.py}
checks the integral polynomial identity and every trace character
exactly. The local torsion proof above already works without this
additional kernel calculation.
\end{proof}
